\providecommand{\kchH}{\kappa_{\mathrm{ch}}^{\mathrm{Heis}}}
\providecommand{\kcat}{\kappa_{\mathrm{cat}}}
\providecommand{\kBKM}{\kappa_{\mathrm{BKM}}}
\providecommand{\kfib}{\kappa_{\mathrm{fiber}}}
\providecommand{\hCS}{\mathrm{hCS}}
\providecommand{\Obs}{\mathrm{Obs}}
\providecommand{\Yaff}{Y_{\epsilon_1,\epsilon_2,\epsilon_3}}
\providecommand{\gghone}{\widehat{\mathfrak{gl}}_1}
\providecommand{\gBKM}{\mathfrak{g}_{\Delta_5}}
\providecommand{\CoHA}{\mathrm{CoHA}}
\providecommand{\Phiten}{\Phi_{10}}
\providecommand{\C}{\mathbb{C}}
\providecommand{\R}{\mathbb{R}}
\providecommand{\cO}{\mathcal{O}}
\providecommand{\cN}{\mathcal{N}}
\providecommand{\cS}{\mathcal{S}}
\providecommand{\AdS}{\mathrm{AdS}}
\providecommand{\BTZ}{\mathrm{BTZ}}
\providecommand{\BH}{\mathrm{BH}}

\section*{Holography and quantum gravity for \(K3\times E\)}
\label{sec:holographic-qg-synthesis}

\paragraph{Synthesis.}
Let \(X=K3\times E\).  The same \(K3\)-fibre Calabi-Yau datum appears
in five comparison problems: class-\(\cS\) compactification, five
dimensional partially holomorphic Chern-Simons boundary algebras, six
dimensional holomorphic Chern-Simons observables, the compact
Hall-Borcherds BPS algebra, and type IIB dyon counting on
\(K3\times T^2\).  The common automorphic invariant is the primitive
Borcherds denominator \(\Delta_5\), with
\[
  \kBKM(\Delta_5)=c_1(0)/2=10/2=5.
\]
The physical dyon denominator is \(\Phiten=\Delta_5^2\).  Thus the
physical scalar partition function has Siegel weight \(10\), while the
Borcherds-Kac-Moody algebra \(\gBKM\) is attached to the primitive
square root of weight \(5\).

\subsection*{Class-\(\cS\) surface data}

\paragraph{Theorem.}
Gaiotto's construction assigns to a punctured curve \(\Sigma\) and type
\(A_1\) six dimensional \((2,0)\) theory a four dimensional
\(\cN=2\) theory \(T[A_1,\Sigma]\), with a weak-coupling frame for each
pants decomposition of \(\Sigma\) (Gaiotto 2009, arXiv:0904.2715).
The Seiberg-Witten curve is the spectral curve in \(T^\ast\Sigma\),
and the charge lattice is read from the corresponding Hitchin system.

\paragraph{K3-fibre comparison.}
For the \(K3\times E\) programme the proposed class-\(\cS\) input is
fixed by matching the Mukai charge lattice, the twenty-four K3
fibre classes, and the automorphic denominator \(\Delta_5\).  The two
natural curves are \(\Sigma_{2,0}\) and \(\Sigma_{0,24}\).  A proof of
their equality or duality requires an explicit Coulomb-branch
identification, including the puncture residues and the induced action
on the \(\Delta_5\) Weyl chamber.  The value \(c_{4d}=107/6\), when
used, is a Shapere-Tachikawa anomaly calculation for this proposed
Gaiotto datum, not a consequence of the Borcherds weight alone.

\subsection*{Five dimensional boundary algebra}

\paragraph{Theorem in the abelian case.}
Partially holomorphic Chern-Simons theory on \(\R\times\C^2\) with
gauge algebra \(\mathfrak{u}(1)\) carries a boundary vertex algebra on
\(\{0\}\times\C\).  In the Costello, Costello-Gaiotto, and
Costello-Dimofte-Paquette framework this boundary algebra is the
affine Yangian of \(\gghone\),
\[
  \partial\hCS_5 \simeq
  \Yaff(\gghone),\qquad
  \epsilon_1+\epsilon_2+\epsilon_3=0.
\]

\paragraph{Compact K3 refinement.}
The \(K3\times E\) refinement replaces the toric local algebra by the
critical Hall positive half attached to \(K3\times E\).  The Borcherds
branch is the Hall-Drinfeld double of that positive half.  It is
separate from the Mukai self-mirror K3 Yangian branch on the rank
\(24\) Mukai lattice.  The compact comparison requires the oriented
\(\hCS\)-to-Hall map, the finite-height Hall-Drinfeld promotion, and
the denominator comparison with \(\Delta_5\).

\subsection*{Six dimensional holomorphic Chern-Simons observables}

\paragraph{Theorem.}
The factorisation algebra \(\Obs_{\hCS}\) of quantum observables of
holomorphic Chern-Simons theory on \(\C^3\) is an \(E_3\)-algebra in
Dolbeault complexes (Costello-Gwilliam, \emph{Factorization Algebras
in Quantum Field Theory}, Vol. II, Theorem 9.3.1; Gwilliam-Williams
2021, Theorem 2.5.5).  The \(E_3\)-product is BV convolution with the
Costello heat-kernel propagator; higher coherences are controlled by
\(\bar\partial\)-flow on the Fulton-MacPherson compactification of
configuration spaces in \(\C^3\).

\paragraph{CY-to-chiral scope.}
The six dimensional hCS algebra is the bulk factorisation object.  After
the \(d=3\) CY-to-chiral Stage-2 specialisation along the K3 fibre, the
chiral output \(A\) on the elliptic curve is \(E_1\).  Braided \(E_2\)
structure appears only after passing to the derived chiral centre or
the Drinfeld centre of \(\mathrm{Rep}^{E_1}(A)\), under the compact Hall
comparison hypotheses.

\subsection*{CoHA, BPS algebra, and the Borcherds denominator}

\paragraph{Theorem on \(\C^3\).}
Schiffmann-Vasserot identify the cohomological Hall algebra of the
tripled \(\C^3\)-quiver with the enveloping algebra of the positive
half of the affine Yangian:
\[
  \CoHA(\C^3)=
  U\!\left(Y^+_{\epsilon_1,\epsilon_2,\epsilon_3}(\gghone)\right).
\]
The full affine Yangian is obtained after the appropriate double or
centre operation, together with the Maulik-Okounkov stable-envelope
identification.  The bare CoHA is the positive half.

\paragraph{Conjectural compact identification.}
For \(X=K3\times E\), the expected compact BPS algebra is the Hall
positive half of the Schiffmann-Vasserot K3 construction with the
elliptic factor included.  Its Hall-Drinfeld double is expected to be
\(\gBKM\), with root multiplicities read from the Borcherds product
for \(\Delta_5\).  The numerical checks are
\[
  \kcat(X)=\chi(\cO_X)=\chi(\cO_{K3})\chi(\cO_E)=2\cdot0=0,
  \qquad
  \kchH(X)=3,
\]
\[
  \kBKM(\Delta_5)=5,
  \qquad
  \kfib(K3)=24.
\]
The set \(\{0,3,5,24\}\) records four distinct constructions:
total-space Euler characteristic, Heisenberg-Cartan chiral rank,
Borcherds weight, and Mukai-lattice fibre rank.

\subsection*{Type IIB dyons and \(\AdS_3\)}

\paragraph{Theorem.}
For type IIB on \(K3\times T^2\), the indexed degeneracy of
quarter-BPS dyons with charge \(Q=(Q_e,Q_m)\) is extracted from the
Igusa cusp form:
\[
  d(Q)=(-1)^{Q\cdot Q+1}
  \oint
  \frac{e^{2\pi i\, Q\cdot\Omega}}{\Phiten(\Omega)}\,d\Omega,
  \qquad
  \Omega\in\mathbb{H}_2.
\]
This is the Dijkgraaf, Verlinde, and Verlinde; Dijkgraaf, Moore,
Verlinde, and Verlinde; Shih, Strominger, and Yin; and Sen formula for
the physical scalar denominator.

\paragraph{Primitive BKM denominator.}
Gritsenko-Nikulin identify \(\Phiten\) with the scalar square
\(\Delta_5^2\) in Igusa normalisation.  The BKM algebra \(\gBKM\) uses
the primitive denominator \(\Delta_5\), whereas the macroscopic dyon
index uses \(1/\Phiten=1/\Delta_5^2\).  The logarithm of the physical
partition function therefore counts both chiral halves.  The
chiral-BKM interpretation of a single boundary algebra is the square
root statement and requires the choice of multiplier and spin cover.

\paragraph{Entropy.}
For large charges the Bekenstein-Hawking entropy is
\[
  S_{\BH}(Q)=2\pi
  \sqrt{Q_e^2Q_m^2-(Q_e\cdot Q_m)^2}.
\]
Sen's quantum entropy function matches the asymptotic expansion of
\(\log d(Q)\) in the large-charge chamber.  The exact chiral boundary
identification with \(\gBKM\) is a BPS-algebra conjecture, not a
macroscopic entropy theorem.

\subsection*{Integrated invariant}

The \(K3\times E\) synthesis is pinned by the Borcherds weight identity
\[
  \kBKM(\Phi_N)=c_N(0)/2.
\]
On the CHL ladder \(N\in\{1,2,3,4,6\}\),
\[
  (c_N(0))=(10,8,6,4,2),\qquad
  \bigl(\kBKM(\Phi_N)\bigr)=(5,4,3,2,1).
\]
At \(N=1\) the primitive denominator is \(\Delta_5\), and the physical
Igusa form is its square.  For \(N>1\), identifying a physical scalar
denominator with the square of a primitive BKM denominator is an
additional normalisation theorem.  No component of the
\(\{0,3,5,24\}\) spectrum is obtained by adding a fixed K3-fibre Euler
term to the BKM weight.

\paragraph{Primary sources.}
Borcherds 1995, arXiv:alg-geom/9508032; Gritsenko-Nikulin 1998,
\emph{International Journal of Mathematics} 9; Gaiotto 2009,
arXiv:0904.2715; Costello 2013, arXiv:1303.2632; Costello-Gwilliam,
\emph{Factorization Algebras in Quantum Field Theory}; Costello-Gaiotto
2018, arXiv:1812.09257; Costello-Dimofte-Paquette 2020,
arXiv:2005.00083; Schiffmann-Vasserot 2013, arXiv:1202.2756;
Maulik-Okounkov 2012, arXiv:1211.1287; Oberdieck-Pixton 2017,
arXiv:1706.10100; Dijkgraaf-Moore-Verlinde-Verlinde 1997,
\emph{Communications in Mathematical Physics} 185; Sen 2008,
arXiv:0708.1270; Shih-Strominger-Yin 2005, arXiv:hep-th/0506151;
Lorgat 2020, arXiv:2004.01676.
